% 本文件由 LaTeX 排版 Agent 根据有效 translation.json 自动生成。
% author=Methodius; work_id=0624; title=论自由意志

\chapter{论自由意志}

% structure: p0001 source=h1 target=omitted reason=page-title-already-rendered-by-parent

\Emph{正统者.} 伊萨卡的老翁，按照希腊人的传说，当他渴望听见塞壬的歌声，因她们那肉欲之声的魅力，便乘船前往西西里，身上绑着绳索，并堵住了同伴们的耳朵；并非他嫉妒他们听见，或愿意给自己加上束缚，而是因为那些歌者的音乐对听见的人结果是死亡。因为在希腊人看来，塞壬的魅力就是如此。如今我并未听到任何这样的歌声；我也不渴望听见那吟唱挽歌的塞壬，她们的沉默比她们的声音对人更有益处；但我祈求享受一种神圣之声的喜乐，这声音虽常被听见，我仍渴望再听；并非因我被肉欲之声的魅力所征服，而是我正在被教导神圣的奥迹，并且期待结果不是死亡，而是永远的救恩。因为歌者并非希腊那致命的塞壬，而是一队神圣的先知，与他们在一起，无需堵住同伴们的耳朵，也无需给自己加上束缚，以免听到的惩罚。因为在一种情形下，听见者随着声音的进入便停止了生存；在另一种情形下，他听见得越多，就越享受更好的生命，被神圣之神引领向前。那么，让每一个人都来，毫无恐惧地聆听这神圣之歌。在我们这里，没有来自西西里海岸的塞壬，没有尤利西斯的束缚，没有融化的蜡灌入人的耳朵；而是一切束缚的解除，以及每一个接近的人都有自由去听。因为这样的一首歌值得我们聆听；而听这样的歌者，在我看来，是一件该祈求的事。但如果有人也愿意听宗徒们的合唱，他会发现同样的和谐之歌。因为前者以奥妙的方式预先唱出了天主的计划；而后者则唱出了对前者神秘宣告的诠释。哦，由圣神谱写的和谐一致！哦，那些歌唱天主奥迹者的优美！哦，但愿我也能在祈祷中参与这些歌唱。那么，让我们也唱起同样的歌，向圣父高举颂赞，在圣神内光荣那在他怀中的耶稣。\BibleRef{若 1:18}\par

人啊，不要躲避一首属灵的诗歌，也不要对听它抱有恶意。死亡不属于它；救恩的故事是我们的歌唱。当我在谈论这些主题时，似乎已经尝到了更美好的享受；尤其当我面前有这样一片花开茂盛的草场，那就是我们这由共同歌唱和聆听神圣奥迹的人组成的聚会。因此，我敢请求你们以毫无嫉恨的耳朵听我说，不要效法加音的嫉妒，\BibleRef{创 4:5} 或像厄撒乌一样迫害自己的兄弟，\BibleRef{创 27:41} 或认同若瑟的兄弟们，\BibleRef{创 37:4} 因为他们因他的话语恨恶自己的兄弟；但要与所有人大不相同，以致你们每个人都惯于说出邻人的心意。为此，你们中间没有邪恶的嫉妒，因为你们已承担起补足弟兄缺欠的责任。啊，尊贵的听众，可敬的团体，属灵的食物！但愿我永远有份于这样的喜乐，这是我的祈祷！\par

\Emph{瓦伦廷.} 我的朋友，昨晚我沿着海边散步，有些专注地凝视大海时，看到了一件神力的非凡实例和由智慧科学产生的工艺品，如果这样的东西可以称为工艺品的话。正如荷马的那句诗所说——\Quote{正如从色雷斯吹来的两股逆风，\newline{}玻瑞阿斯和仄费罗斯，搅动鱼深的海，\newline{}突然，四周，黑色的洪水\newline{}高高卷起，把盐草抛到岸上；}——所以在我看来昨晚的事正是如此。因为我看见浪涛如山巅，几乎可以说是高耸入天。因此我原本只期待整个大地被淹没，并开始在心中构建一个避难之处，以及一个诺厄方舟。但事情并非如我所想；因为，正当海浪涨到顶峰时，它又自行消散，没有越过自己的界限，仿佛对天主的法令存有敬畏。\BibleRef{约 38:11} 正如一个仆人常常被迫去做违背自己意愿的事，由于惧怕而服从命令，不敢对在不愿做事时的痛苦说一句话，却满腔愤怒地喃喃自语——在我看来，大海似乎也是这样，仿佛发怒却又将敬畏限制在自己之内，抑制着自己，不让它的主人察觉它的愤怒。\par

在这些现象中，我开始沉默地凝视，并想在心里测量天及其球体。我开始探究它从哪里升起，在哪里落下；以及它有怎样的运动——是前行运动，即从一个地方到另一个地方，还是旋转运动；此外，它的运动是如何持续的。并且，诚然，也值得探究关于太阳——它被安置在天中的方式；还有它运行的轨道；以及它短暂时间后去向何处；为什么它甚至也不偏离自己的轨道：但他，如同人所说，也在遵守更高权能者的诫命，只在被允许的时候才与我们一同显现，然后离开，仿佛被召唤一般。\par

因此，当我探究这些事情时，我看见阳光正在消退，白昼渐渐消逝，黑暗立刻降临；太阳之后是月亮，她在初升时并不圆满，但在运行中渐渐显得更大。我也没有停止探究她，而是考察了她亏损和圆满的原因，以及为什么她也遵守日子的轮转；从这一切看来，似乎有一种神圣的管理和权能支配着整个宇宙，我们按正义可以称之为天主。\par

于是我开始赞美造物主，因为我看见大地稳固不动，生物种类繁多，植物花朵色彩斑斓。但我的心思并不只停留于这些，而是开始探究它们从何而来——是来自某个与天主永恒共存的源头，还是唯独来自祂自己，没有与祂共存者？因为我认为，正确看法是祂没有从无中制造任何东西，除非我的理性完全不可靠。因为生成之物，其本性是从已有之物衍生起源。而且在我看来，同样可以真实地说，没有东西与天主永恒共存而不同于祂自身，但凡存在之物都源于祂，这一点我也被元素的无可否认的布置和自然围绕它们的井然秩序所说服。\par

于是，带着一些关于事物美好秩序的想法，我回到了家。但在次日，也就是今天，当我来时，我看见两个同类——我指人——在互相殴打辱骂；另一个人则想剥去邻居的衣物。现在有些人开始尝试更可怕的行为：一个人剥去一具尸体的衣服，将一度埋藏在地下的身体再次暴露于日光之下，并且对待和自己相似的形象如此侮辱，以至将尸体留给狗作食物；另一个人拔出剑，攻击一个和自己一样的人。他想通过逃跑获得安全，但对方却不停止追赶，也不控制愤怒。我还需要多说吗？他攻击他，立刻用剑砍伤了他。\Emph{因此那位受伤的人}向他的同类哀告，伸出双手恳求，愿意放弃自己的衣服，只求活命。但对方没有平息愤怒，也没有怜悯同胞，不肯在他面前的这存在中看到自己的形象；反而像一头野兽，用剑准备吞食他。现在他甚至把嘴凑到那和自己如此相似的身体上，其狂怒之甚，至于此。可以看到一个人遭受虐待，另一个人立刻剥去他的衣服，甚至不将脱去衣服的身体埋入土中。除此之外，还有另一个人，抢夺他人婚姻的权利，想要侮辱邻居的妻子，怂恿她转向非法的拥抱，不希望她的丈夫成为自己孩子的父亲。\par

此后，我开始相信悲剧，认为\Person{提厄斯忒斯}的宴席真的发生过；也相信\Person{奥诺玛奥斯}的非法情欲，也不怀疑兄弟拔剑相向的争斗。\par

因此，在目睹了这样的事情之后，我开始探究它们从何而来，它们的起源是什么，谁是对人类设下如此诡计的始作俑者，它们从何被发明，谁是传授者。若敢说天主是这些事情的作者，是不可能的；因为甚至不能说它们从祂获得其实体或存在。因为怎么能对天主怀有这些想法呢？祂是美善的，是卓越之物的创造者，祂没有任何邪恶。不，祂的本性是不以此为乐；祂禁止它们的产生，拒绝那些以此为乐的人，但接纳那些避免它们的人。若称天主为祂不认可之物的制造者，岂不荒谬？因为如果祂首先创造了它们，祂就不会希望它们不存在；而且祂希望接近祂的人效法祂。\par

因此，在我看来，将这些事归于天主，或说它们从祂而出，是不合理的；尽管必须承认，如果祂制造了邪恶之物，那么从无中产生某物是可能的。因为那从无中带来存在者，不会使它们失去存在。再者，必须说曾有一个时候天主以邪恶之事为乐，而现在则不是这样。因此，在我看来，不可能对天主这样说。因为将这一点归给祂，不适合祂的本性。因此，在我看来，与祂共存着某种称为物质的东西，祂从中形成了现存之物，用智慧的艺术加以区分，并以美丽的秩序安排它们，邪恶之物似乎也是由此而生。因为这种物质没有品质和形式，此外还无序地飘荡，未被神圣的艺术触及，天主不怨恨它，也不让它继续如此飘荡，而是开始塑造它，并想将其最好的部分与最坏的部分分开，从而制造出所有适合天主用它制造的东西；但其中像渣滓的部分，可以说，不适合制造任何东西，祂便随它去，因为它对祂没有用处；由此，在我看来，邪恶之物现在流入了人间。这在我看来是对这些事的正确看法。但是，我的朋友，如果你认为我说了什么错误的话，请指出，因为我非常渴望了解这些事。\par

\Person{正统者}。我欣赏你的直率，我的朋友，并赞赏你对这个主题的热忱；至于你关于现存之物所表达的意见，即天主是从某种底层基质制造了它们，我并非完全不同意。因为，的确，恶的起源\Emph{是一个主题}，曾引出许多人的见解。在你我之前，无疑有许多有才之士对这个问题做过最透彻的探究。其中一些人表达了和你相同的意见，但另一些人则把天主说成这些事物的创造者，害怕承认有物质与祂同时存在；而前者，因为害怕说天主是邪恶的作者，认为应当把物质说成与祂同时存在。这两者都未能正确谈论这个问题，原因在于他们对天主的敬畏与对真理的准确知识不相符合。\par

但另一些人完全拒绝对此问题加以探究，理由是这样探究是无穷无尽的。至于我，因着与你之间的友谊，不允许我推辞探讨的主题，尤其当你宣称你的意向时——你并非受偏见所左右，尽管你曾凭推测对事物的状况持有见解——而是说自己坚定了寻求真理的渴望。\par

因此，我将欣然转向对这问题的讨论。但我希望我身边的这位同伴也聆听我们的对话。因为实际上，他在这些事上似乎与你持有很多相同的见解，所以我希望你们俩都能参与讨论。因为无论我向身处此境的你说什么，我也同样会对他说。所以，如果你足够宽容，认为我在这个重大主题上说得正确，就请回答我提出的每一个问题；因为其结果将是，你获知真理，而我也不会随意与你进行讨论。\par

\Person{瓦伦提尼安}：我乐意照你说的做；所以请准备好提出那些你认为我能从中准确了解这个重要主题的问题。因为我为自己设定的目标并非卑劣的取胜，而是彻底认识真理。所以请继续讨论吧。\par

\Person{正统者}：好，那么，我想你不会不知道，两个非受造的事物共存是不可能的，尽管你在谈话的前一部分似乎几乎表达了这一点。的确，我们必然要说两者之一：要么天主与物质分离，要么，反之，祂与物质不可分离。那么，若有人说它们联合，他便要说非受造的只有一位，因为所谈的每一物都将是另一物的一部分；并且由于它们互为部分，便不是两个非受造者，而是由不同元素组成的一个。因为我们不会因为一个人有不同肢体，就把他分裂成多个存在。而是按照理性的要求，我们说一个由许多部分组成的单一存在——人——是由天主创造的。所以，如果天主不与物质分离，就必须说非受造的只有一位；但若有人说祂是分离的，那么两者之间必然有某个中介物，使它们的分离显明。因为除非有某物作为比较距离的参照，否则无法估算一物与另一物的距离。这不仅适用于我们眼前的实例，也适用于其他任何数目。因为我们在两个非受造物情况下提出的论证，若承认非受造物有三个，也必然具有同等效力。因为关于它们，我也会问：它们是彼此分离，还是反之，彼此联合？因为若有人决意说它们联合，他将被告知如前所述；若又说它们分离，他也无法避免那分离者的必然存在。\par

那么，若有人说关于非受造物还可以有第三种恰当的说法，即：天主既不与物质分离，也不作为整体的一部分与物质联合，而是天主局部地位于物质之中，物质又在天主之中，那么，作为结果，必须告诉他：若我们说天主被置于物质之中，我们就必然说祂受到限制，并被物质所包围。但随后祂必然与物质一样无序地被运载。并且祂不安息，也不自在，这是祂在其中被运载的必然结果——时而朝此，时而朝彼。除此之外，我们还必须说天主的处境更糟。\par

因为若物质一度无序，而祂决定使之变好并加以整理，那么曾有一个时候，天主在那无序之中。我也可以合情合理地问这个问题：天主是充满物质整体，还是存在于物质的某部分？因为若有人决意说天主在物质的某部分，那么他将天主置于比物质小得多的境地——也就是说，如果物质的一部分完全容纳了天主。但若他说祂在物质的所有部分中，并遍及整个物质，他必须告诉我们祂是如何作用于物质的。因为我们必须说，天主有一种收缩，借此收缩，祂作用于祂所撤回的那部分，或者祂与物质联合作用，没有撤回的余地。但若有人说物质在天主内，同样需要探究：这是否由于祂自我分离，如同受造物存在于空气中，祂被分割并分开以容纳其中的存在？还是物质局部地位于祂内，即如水在土中？因为若我们说如同在空气中，我们必须说天主是可分的；但若如同水在土中——既然物质无序且无安排，此外还含有恶——我们必须说，在天主内包含无序和恶。在我看来，这是一个不妥当的结论，不，甚至是一个危险的结论。因为你希望物质存在，以避免说天主是恶的创造者；而决定避免这一点，你却说祂是恶的容器。\par

那么，假设你认为物质与受造实体分离，并说它是非受造的，我本应对此大加谈论，以证明它不可能非受造；但既然你说恶的起源问题是这一假设的原因，因此我觉得正确的方法是进而探究这一点。因为当恶如何存在这一点被清楚说明，并且不可能说天主是恶的原因——由于物质从属于祂——那么，在我看来，要推翻这样的假设，只需指出：若天主创造了之前不存在的属性，祂同样也创造了实体。那么，你是否说，存在与天主同在的无属性物质，祂从中形成了这个世界的开始呢？\par

\Person{瓦伦提尼安}：我是这么想的。\par

\Person{正统者}：那么，如果物质没有性质，世界是由天主创造的，而世界中存在着性质，那么天主就是性质的创造者吗？\par

\Person{瓦伦提尼安}：是这样。\par

\Person{正统者}：既然我先前听你说过，任何事物都不可能从不存在的东西产生，那么请回答：你是否认为世界的性质并非由任何现成的性质产生？\par

\Person{瓦伦提尼安}：我是这么认为的。\par

\Person{正统者}：而且它们不同于实体？\Person{瓦伦提尼安}：是的。\par

\Person{正统者}：那么，如果性质既不是天主从任何现成的东西中造的，也不是从实体中获得存在（因为它们不是实体），我们就必须说，它们是由天主从不存在的东西中造的。因此，我认为你先前说不可能设想任何事物由天主从不存在的东西中产生，那是言过其实了。\par

不过，关于这个问题的讨论就到此为止吧。因为我们确实见到，在我们中间，有些人从不存在的东西中造出东西来，尽管他们在很大程度上似乎是用某些东西来制造。例如，我们可以从建筑师的情况中找到例子：他们确实不是从城市中造出城市，也不是从神庙中造出神庙。\par

但是，如果你因为这些东西下面有实体，就认为建造者是从存在的东西中造出它们，那你就计算错了。因为使城市或神庙成形的并不是实体，而是应用于实体的技艺。而这种技艺并非从存在于实体本身的某种技艺中产生，而是从并不存在于实体中的东西中产生。\par

不过，你似乎可能用这样的论点来反驳我：工匠是从他所拥有的技艺中造出与实体相联的技艺。对此，我认为一个恰当的回答是：在人身上，技艺并非从任何潜藏于下的技艺中产生；因为不能承认实体本身就是技艺。技艺属于偶性之类，并且是只有当它们被用于某种实体时才存在的东西之一。因为即使没有建造的技艺，人也会存在；但除非人先存在，否则技艺就不会存在。由此我们必须说，按事物的本性，技艺是从不存在的东西中在人身上产生的。那么，既然我们已经表明在人身上是这样，为什么说天主能够不仅从不存在的东西中造出性质，而且造出实体，就是不恰当的呢？因为既然某种东西从不存在的东西中产生似乎是可能的，那么实体的情况显然也是如此。现在回到恶的问题。你认为恶属于实体之类，还是属于实体的性质？\par

\Person{瓦伦提尼安}：属于性质。\par

\Person{正统者}：但我们发现物质是没有性质或形式的？\par

\Person{瓦伦提尼安}：是的。\par

\Person{正统者}：那么，这些名称与实体的关联是由于实体的偶性。因为谋杀不是实体，其他任何恶也不是实体；但实体因实行谋杀而获得了一个同名的称呼。因为一个人（不被说成是）谋杀，而是通过犯谋杀而获得\Quote{凶手}这个派生的名称，但他本身并不是谋杀；简言之，其他任何恶都不是实体；但由于实行某种恶，它可以被称为恶。同样，如果你认为有别的什么东西是人受恶的原因，那么它也是因为它通过人起作用并唆使人犯恶，才被称为恶。因为人因自己的行为而成为恶的。他被称为恶，是因为他是恶的行事者。一个人所做的，不是他本人，而是他的活动；正是从他的行为中，他获得了恶的名号。因为如果我们说，他就是他所做的，而他犯了谋杀、通奸等诸如此类的事，那么他就是所有这些。如果他是这些，那么当这些产生时，他就存在；但当这些不存在时，他也就不存在了。而这些是由人产生的。那么，人就是它们的作者，是它们存在或不存在的因由。但如果每个人因自己所行的事而成为恶的，而他所行的事是有起源的，那么他在恶中也有一个开端，而恶也有一个开端。既然如此，就没有人在恶中没有开端，恶的事物也没有起源。\par

\Person{瓦伦提尼安}：好吧，我的朋友，在我看来，你似乎已经很充分地论证了反驳对方的观点。因为你从我们在讨论中同意的前提中得出了正确的结论。的确，如果物质是没有性质的，那么天主就是性质的创造者；而如果恶是性质，天主就是恶的作者。但在我看来，说物质没有性质是假的；因为不能说任何实体是没有性质的。实际上，就在你说它没有性质这一点上，你就是在宣称它具有一种性质，因为你描述了物质的特点，而这是一种性质。所以，如果你愿意，请从头开始讨论；因为在我看来，物质从未开始拥有性质。因为情况既然如此，我的朋友，我断言恶是从物质的流溢中产生的。\par

\Person{正统者}：如果物质从永远就拥有性质，那么天主将是什么的创造者呢？如果我们说是实体，那么我们就认为它们是先存的；如果我们又说是性质，那么它们也被宣称为存在。既然实体和性质都存在，在我看来，称天主为创造者就是多余的。但请回答一个问题：你说天主是创造者，是借着将那些实体的存在改变为不存在，还是借着改变性质而保存实体？\par

\Person{瓦伦提尼安}：我认为并没有本体的变化，只有属性的变化；就属性而言，我们称天主为创造者。正如有人可能会说房屋是由石头建造的，但并不能说它们不再是石头本体，因为它们被称为房屋；因为我认为房屋是由建造的性质构成的。同样，我认为天主使本体保持不变，而使其属性发生变化，正因如此，我说这个世界是由天主创造的。\par

\Person{正统者}：你是否也认为恶属于本体的属性之一？\par

\Person{瓦伦提尼安}：我认为是的。\par

\Person{正统者}：这些属性是原初就存在于物质中，还是有其起始？\par

\Person{瓦伦提尼安}：我认为这些属性与物质永恒共存。\par

\Person{正统者}：但你难道不是说天主使属性发生了变化吗？\par

\Person{瓦伦提尼安}：我是这么说的。\par

\Person{正统者}：是变得更好吗？\par

\Person{瓦伦提尼安}：我认为是这样。\par

\Person{正统者}：那么，如果恶属于物质的属性，而天主使这些属性变得更好，就必须追问恶是从何而来的。因为，要么所有恶的属性都变得更好，要么有些是恶的，有些不是，那恶的就没有变得更好；但其余的，就其所具有的优越性而言，被天主为秩序的缘故而改变了。\par

\Person{瓦伦提尼安}：这正是我从一开始就持有的观点。\par

\Person{正统者}：那么，你怎么说天主让恶的属性保持原样？是因为祂能除去却不愿意，还是愿意却无能为力？因为如果你说祂能却不愿，那祂就是这些恶的制造者；因为祂虽有能力消除恶，却任其保持原样，尤其是在祂开始作用于物质之后。因为如果祂与物质毫无关系，那么祂就不是所容许之事的制造者。但既然祂作用于一部分，而让另一部分保持原状，同时祂又有能力使其变好，我认为祂就是恶的制造者，因为祂让部分物质留在卑劣之中。那么，祂建造一部分是为了毁坏另一部分；从这点来看，我觉得这一部分因祂在物质中的安排而主要受到损害，因为它分享了恶。因为在物质被整理之前，它没有对恶的感知；但现在它的每一部分都有了感知恶的能力。以人为例：在成为活物之前，他对恶没有感觉；但从他被天主塑造成人形之时起，他就获得了对临近恶的感知。所以，你认为天主为物质益处而做的这个行动，结果反而对物质更糟。但如果你说天主不能阻止恶，这不可能是因为祂本性软弱，还是因为被恐惧征服，服于某个更强大的存在？看看你想把哪一个归于全能而良善的天主？但再请回答我关于物质的问题：物质是单纯的还是复合的？因为如果物质是单纯而单一的，而宇宙是复合的，由不同的本体构成，就不能说宇宙是由物质构成的，因为复合物不能由一种纯粹的单一成分组成。组成表明多种单纯物的混合。另一方面，如果你说物质是复合的，那它完全由单纯元素组成，它们曾各自单独单纯，通过它们的组合产生了物质；因为复合物从单纯物获得其组成。所以曾有一个时候物质不存在——即在单纯元素结合之前。但如果曾有一个时候物质不存在，而受造物从没有一个时候不存在，那么物质就不是受造的。由此推出有许多非受造之物。因为如果天主是非受造的，而组成物质的单纯元素也是非受造的，那么非受造物的数目就超过两个。但暂且不论什么是单纯元素、物质或形式——因为这会引出许多荒谬——我问你：你认为存在之物没有任何东西与自己相悖吗？\par

\Person{瓦伦提尼安}：我认为是的。\par

\Person{正统者}：但水与火相悖，黑暗与光明相悖，热与冷相悖，湿与干相悖。\par

\Person{瓦伦提尼安}：我认为是这样。\par

\Person{正统者}：那么，如果没有任何存在之物与自己相悖，而这些彼此相悖，它们就不会是同一个物质——也不会由同一个物质形成。但再问：你认为事物的各部分不会彼此毁灭吗？\par

\Person{瓦伦提尼安}：我认为是的。\par

\Person{正统者}：那么，火、水等也是物质的各部分吗？\par

\Person{瓦伦提尼安}：我认为是的。\par

\Person{正统者}：那么，你难道不认为水会毁灭火，光明毁灭黑暗，等等？\Person{瓦伦提尼安}：我认为是的。\par

\Person{正统者}：那么，如果事物的各部分不会彼此毁灭，而这些被发现会，它们就不是同一事物的部分。但如果不是同一事物的部分，它们就不是同一个物质的各部分。事实上，它们也不是物质，因为没有任何存在之物会自我毁灭。既然对立之物如此，就表明它们不是物质。关于物质，这就够了。\par

现在我们须考察恶事，并必须探究人间的恶。关于这些恶，它们是恶的原则的形式，还是其部分？如果是形式，恶就没有与它们分离的独立存在，因为形式中寻求种类，且潜存于形式之下。但若是如此，恶就有起源。因为它的形式显露出有起源——诸如谋杀、奸淫等。但若你愿它们为某一恶的原则的部分，且它们有起源，则恶的原则也必有起源。因为那些部分有起源的东西，本身也必然是受造的。因为整体由部分构成。若部分不存在，整体也不存在，尽管某些部分可能存在，即使整体不在。\par

现在，不存在某物一部分受造而另一部分不受造的情况。但即使我准许这点，也曾有一段时间恶尚未完备，即在天主塑造物质之前。而当人由天主产生时，恶便达到完备；因为人是恶的部分的制造者。由此推之，使恶完备的原因是天主创造者，这是不敬之说。但若你说恶既不是上述情况，而是某种恶的行为，则你宣称它有起源。因为行为使存在开始。此外，你无法再指出别的恶。因为除了人间发生的事，你还能指出何种其他行为？现已证明，行动者按其存有并非恶，而是按其恶行成为恶。\par

因为没有什么是本性为恶的，而是因使用而成为恶。因此，他说，我说人被造时具有自由意志，并非因为已有恶存在而他可随意选择，而是因为他有服从或不服从天主的能力。\par

因为这就是自由意志恩赐的意义。人在受造后从天主领受一条诫命；由此立即兴起恶，因为他未遵守天主的命令；而这唯独是恶，即不服从，它有开端。\par

因为人领受了能力，却奴役了自己，并非因其本性的不可抗拒倾向所压倒，也非所获的能力剥夺了他更好的东西；因为正是为此，我说他被赋予了这能力（他领受了上述能力），好使他能获得对他已有之物的增补，这增补因他的服从而来自至高的存在，并作为债务向他的造物主索取。因为我说人被造不是为了毁灭，而是为了更好的事物。因为若他被造得像任何元素或那些为天主提供类似服务的事物，他便不再因审慎抉择而配得赏报，而会像造物主的工具；那么他因恶行遭受责备就不合理，因为真正的作者是使用他的那位。但人不理解更好的事物，因为他不知道（他存在的）作者，只知道他被造的目的。因此我说，天主既计划如此尊荣人，并赐予他理解更好事物的能力，便给了他做他所愿之事的能力，并赞扬他运用这能力于更好的事物；并非祂再次剥夺他的自由意志，而是愿指出更好的道路。因为能力与他同在，他领受诫命；但天主劝勉他将选择的能力转向更好的事物。因为正如父亲劝勉有能力学习功课的儿子要更加注意，因为父亲指出这是更好的道路，却未剥夺儿子的能力，即使儿子不愿乐意学习；同样，我认为天主敦促人服从祂的命令时，并未剥夺他决意和拒服的能力。因为祂指出给予这劝告的原因，在于未剥夺他的能力。但祂颁布命令，为使人能享受更好的事物。因为这就是服从天主命令的结果。因此，祂颁布命令并非要剥夺祂已赐予的能力，而是为使更好的恩赐得以赐下，如同赐给配得更大事物的人，以回报他在有能力拒绝时却服从了天主。我说人被造具有自由意志，并非因为已有某种恶存在而他可随意选择……而是服从和不服从天主的能力是唯一原因。\par

因为这原是由自由意志所要达成的目的。人在受造后从天主领受一条诫命，由此立即兴起恶；因为他未遵守天主的命令，而这唯独是恶，即不服从，它有开端。因为无人有权说它无起源，既然其作者有起源。但你必会问这不服从从何而来。在圣经中清楚记载，由此我能说人并非由天主造为这种状态，而是通过某种教导成为这样。因为人并未领受这样的本性。因为若其本性如此，这事就不会通过教导临到他。如今在圣经中说：\Quote{人学会了（恶）。}\BibleRef{耶 13:23} 那么，我说不服从天主是被教导的。因为唯独此恶是违反天主旨意而产生的，因为人不会自行学会恶。因此，教导恶的便是那蛇。\par

至于我，我说恶的开端是嫉妒，它源于人因天主赐予更高尊荣而与众不同。如今恶乃不服从天主的命令。\par

\SourceRecord{Translated by William R. Clark.；Ante-Nicene Fathers；Vol. 6.；Edited by Alexander Roberts, James Donaldson, and A. Cleveland Coxe.；Buffalo, NY: Christian Literature Publishing Co.,；1886.；<http://www.newadvent.org/fathers/0624.htm>.}
